% purpose: 说明两相邻时段的购电量在供需平衡、储电量上下界与单段功率上限交出的凸可行域上，
%          线性费用的等值线把最优解推到唯一顶点
% topology: geometry   reading_direction: coordinate-based
% node_roles: 约束边界、可行域、等费用线、最优顶点
% visual_encoding: 线型区分约束边界与等费用线，斜线填充标记可行域，实心点标记最优顶点
% style_family: restrained Nature-style line art   print_mode: colorblind-safe   final_width: text
\documentclass[border=3pt,11pt]{standalone}
\usepackage{ctex}
\usepackage{amsmath}
\usepackage{tikz}
\usetikzlibrary{arrows.meta,patterns}
\definecolor{mhblue}{HTML}{0072B2}
\begin{document}
\begin{tikzpicture}[
  bound/.style={draw=mhblue,thin,dash pattern=on 3pt off 2pt},
  region/.style={draw=mhblue,line width=0.9pt},
  iso/.style={draw=mhblue,thin,dash pattern=on 1pt off 1.6pt},
  axis/.style={draw=mhblue,line width=0.7pt,-{Stealth[length=4pt]}},
  lead/.style={draw=mhblue,line width=0.5pt},
  lb/.style={font=\small,inner sep=1.2pt}]

% ---- 坐标轴 ----
\draw[axis] (0,0) -- (7.6,0);
\draw[axis] (0,0) -- (0,6.3);
\node[lb] at (7.55,-0.45) {\(x_1\)};
\node[lb] at (-0.52,6.25) {\(x_2\)};

% ---- 约束边界：两段功率上限、供需平衡给出的购电下限、储电量上下界 ----
\draw[bound] (1,0) -- (1,6.0);
\draw[bound] (6,0) -- (6,4.2);
\draw[bound] (0,0.8) -- (7.2,0.8);
\draw[bound] (0,5.2) -- (5.2,5.2);
\draw[bound] (0,3.8) -- (3.8,0);
\draw[bound] (3.2,6.0) -- (6.8,2.4);

% ---- 可行域（六边形）----
\fill[pattern=north east lines,pattern color=mhblue]
  (1,2.8) -- (1,5.2) -- (4,5.2) -- (6,3.2) -- (6,0.8) -- (3,0.8) -- cycle;
\draw[region]
  (1,2.8) -- (1,5.2) -- (4,5.2) -- (6,3.2) -- (6,0.8) -- (3,0.8) -- cycle;

% ---- 等费用线与费用下降方向 ----
\draw[iso] (2.2,4.2) -- (5.4,2.2);
\draw[iso] (3.4,4.6) -- (6.6,2.6);
\draw[lead,-{Stealth[length=5pt]}] (4.4,2.2) -- (3.78,1.21);

% ---- 最优顶点 ----
\fill[black] (3,0.8) circle (2.6pt);

% ---- 标注 ----
\node[lb] at (1,-0.45) {\(x_1^{\min}\)};
\node[lb] at (6,-0.45) {\(x_1^{\max}\)};
\node[lb] at (-0.62,0.8) {\(x_2^{\min}\)};
\node[lb] at (-0.62,5.2) {\(x_2^{\max}\)};
\node[lb,fill=white,anchor=west] at (1.15,1.7) {储电量下界};
\node[lb,fill=white,anchor=west] at (6.45,4.9) {储电量上界};
\draw[lead] (6.4,4.78) -- (5.6,4.0);
\node[lb,fill=white] at (3.9,3.05) {可行域};
\node[lb,fill=white,anchor=east] at (2.05,4.35) {等费用线};
\node[lb,fill=white,anchor=west] at (4.55,1.95) {费用下降方向};
\node[lb,fill=white,anchor=west] at (3.35,0.32) {最优顶点};

\end{tikzpicture}
\end{document}
